\section*{Розділ IV. Керівні органи та посадові особи СР СМ}
\addcontentsline{toc}{section}{Розділ IV. Керівні органи та посадові особи СР СМ}

\subsection*{4.1. Голова Студентської ради студмістечка}
\addcontentsline{toc}{subsection}{4.1. Голова Студентської ради студмістечка}
    4.1.1. Студентську раду студмістечка очолює Голова СР СМ, який обирається мешканцями гуртожитків студмістечка на прямих виборах відповідно до Розділу II цього Положення та Положення про ЦВКс.

    4.1.2. Голова СР СМ:

        \begin{enumerate}[label=\alph*)]
            \item Організовує роботу СР СМ та головує на її засіданнях.
            \item Представляє СР СМ у відносинах з адміністрацією Університету та студмістечка, комунальними службами, іншими ОСС, організаціями та установами з побутових питань.
            \item Вносить на розгляд СР СМ пропозиції щодо кандидатур Заступника(-ів) Голови, Секретаря СР СМ та керівників постійних комітетів.
            \item Координує роботу Заступника(-ів) Голови, Секретаря та керівників комітетів СР СМ.
            \item Підписує рішення, протоколи та інші документи СР СМ.
            \item Здійснює оперативне управління коштами та майном, що перебувають у віданні СР СМ, підписує фінансові документи та затверджує витрати в межах кошторису, погодженого СР СМ, та відповідно до її рішень. Це повноваження не поширюється на кошти та майно СРГ.
            \item Забезпечує виконання рішень КСУ та СР СМ.
            \item Звітує про свою діяльність та діяльність СР СМ перед КСУ та СР СМ.
            \item Виконує інші повноваження, передбачені цим Положенням та Положенням про ОСС.
        \end{enumerate}

    4.1.3. Голова СР СМ є підзвітним Конференції студентів Університету. Повноваження Голови СР СМ можуть бути достроково припинені за рішенням КСУ або за рішенням Студентської уповноваженої делегації (СУД) у випадках та порядку, передбачених Положенням про СУД.

    \subsubsection*{4.1.4. Призначення виконувача обов'язків Голови СР СМ}
        4.1.4.1. У разі дострокового припинення повноважень Голови СР СМ, його тимчасової неможливості виконувати свої обов'язки (через хворобу, відрядження тощо) або інших обставин, що унеможливлюють виконання ним своїх повноважень, призначається виконувач обов'язків (в.о.) Голови СР СМ.

        4.1.4.2. В.о. Голови СР СМ може бути призначений колегіальним рішенням СР СМ на термін не більше 1 (одного) місяця. Для призначення в.о. на більший термін необхідне рішення КСУ або СУД.

        4.1.4.3. В.о. Голови СР СМ може бути Голова СРГ або інший студент Університету, що відповідає загальним вимогам до членів ОСС. На період виконання обов'язків Голови СР СМ, Голова СРГ тимчасово делегує свої повноваження заступнику в СРГ без дострокового припинення своїх повноважень у відповідній СРГ. Призначення в.о. Голови СР СМ не вважається обійманням посади Голови СР СМ у розумінні п. 2.3.2 цього Положення.

        4.1.4.4. В.о. Голови СР СМ має повний обсяг повноважень Голови СР СМ, окрім права голосу. Призначення на посаду в.о. Голови СР СМ не надає права голосу в СР СМ. Наявні права голосу, набуті на інших підставах (зокрема, як Голови СРГ відповідно до п. 2.2.1), зберігаються. Якщо рішення про призначення в.о. прийнято КСУ або СУД, у ньому можуть зазначатися додаткові обмеження повноважень.

        4.1.4.5. Передача справ в.о. Голови СР СМ здійснюється відповідно до процедури, визначеної розділом 4.6 Положення про ОСС.

\subsection*{4.2. Заступник(-и) Голови та Секретар СР СМ}
\addcontentsline{toc}{subsection}{4.2. Заступник(-и) Голови та Секретар СР СМ}
    4.2.1. Голова СР СМ може мати одного або кількох заступників та Секретаря СР СМ.

    4.2.2. Заступник(-и) Голови та Секретар СР СМ призначаються Головою СР СМ з числа студентів, що беруть участь у роботі СР СМ (з правом голосу або дорадчим), за погодженням (затвердженням) СР СМ.

    4.2.3. Заступник Голови СР СМ виконує обов'язки Голови СР СМ за його відсутності або за його дорученням, координує роботу визначених комітетів або напрямів діяльності.

    4.2.4. Секретар СР СМ відповідає за ведення протоколів засідань СР СМ, організацію документообігу та архіву СР СМ.

    4.2.5. Повноваження Заступника(-ів) Голови та Секретаря СР СМ можуть бути достроково припинені за рішенням СР СМ.

\subsection*{4.3. Комітети СР СМ}
\addcontentsline{toc}{subsection}{4.3. Комітети СР СМ}
    4.3.1. Для реалізації основних завдань та напрямів діяльності СР СМ може створювати постійні або тимчасові комітети (наприклад: побутовий, безпеки, благоустрою території, культурно-масовий, інформаційний, з питань поселення тощо).

    4.3.2. Структура, перелік та повноваження комітетів затверджуються СР СМ.

    4.3.3. Керівники комітетів призначаються Головою СР СМ з числа студентів, що беруть участь у роботі СР СМ, за погодженням (затвердженням) СР СМ.

    4.3.4. Керівник комітету організовує роботу відповідного підрозділу, звітує про його діяльність перед СР СМ та несе відповідальність за виконання покладених на підрозділ завдань.

    4.3.5. Повноваження керівника комітету можуть бути достроково припинені за рішенням СР СМ.

\subsection*{4.4. Відповідальність посадових осіб}
\addcontentsline{toc}{subsection}{4.4. Відповідальність посадових осіб}
    4.4.1. Заступник(-и) Голови, Секретар СР СМ та керівники комітетів несуть відповідальність перед СР СМ за належне виконання своїх обов'язків.

    4.4.2. У разі систематичного невиконання обов'язків, порушення цього Положення, Положення про ОСС або вчинення дій, що шкодять репутації ОСС, зазначені посадові особи можуть бути достроково відкликані зі своїх посад за рішенням СР СМ.

    \subsubsection*{4.4.3. Призначення виконувачів обов'язків інших посадових осіб СР СМ}
        4.4.3.1. У разі дострокового припинення повноважень, тимчасової неможливості виконувати обов'язки або інших обставин, що унеможливлюють виконання повноважень Заступника(-ів) Голови, Секретаря СР СМ або керівників комітетів, призначення виконувача обов'язків здійснюється в тому ж порядку, що й призначення відповідних посадових осіб.

        4.4.3.2. Термін повноважень, обсяг повноважень та порядок передачі справ для таких в.о. визначаються рішенням про їх призначення.
